\chapter{Peripheral Blood Smear}

\section*{What Is This Test?}
A peripheral blood smear (PBS), also called a peripheral blood film or blood film examination, is a laboratory procedure in which a drop of blood is spread on a glass microscope slide, fixed, stained with Romanowsky-type stains (typically Wright-Giemsa stain), and examined under a microscope by a trained medical laboratory scientist, pathologist, or hematologist. The peripheral blood smear provides direct visual information about the morphology, number, and distribution of all blood cell lineages --- red blood cells (erythrocytes), white blood cells (leukocytes), and platelets (thrombocytes). Unlike automated hematology analyzers, which provide quantitative cell counts and calculated indices, the blood smear reveals qualitative cellular details that are essential for the diagnosis of many hematologic disorders. These include abnormalities of red cell size (anisocytosis), shape (poikilocytosis), color (hypochromia, polychromasia), inclusions (Howell-Jolly bodies, basophilic stippling, Pappenheimer bodies, malaria parasites), white cell morphology (toxic granulation, D\"{o}hle bodies, blast cells, Auer rods), and platelet morphology (giant platelets, platelet clumps). The peripheral blood smear is often triggered by abnormal automated CBC results (flagged by the analyzer) or specific clinical indications. It remains the gold standard for morphological diagnosis of many hematologic conditions and is often the critical test that leads to definitive diagnosis of leukemia, hemolytic anemia, and parasitic infections.

\section*{Alternative Names}
PBS; Peripheral Blood Film; Blood Film; Blood Smear; Blood Slide; Peripheral Smear; Wright-Stained Smear; Manual Differential Smear

\section*{Principle}
A thin film of anticoagulated whole blood (EDTA) is prepared by the wedge technique: a drop of blood approximately 2--3 mm in diameter is placed near one end of a clean glass microscope slide. A second (spreader) slide is placed at a 30--45 degree angle just in front of the drop, backed into the drop so the blood spreads along the edge, then pushed forward in a smooth, rapid motion to create a feathered edge. The slide is air-dried, fixed in absolute methanol, and stained with a Romanowsky-type stain (Wright stain, Wright-Giemsa stain, or May-Gr\"{u}nwald-Giemsa stain). Romanowsky stains contain methylene blue (a basic dye that stains acidic/nucleic acid structures blue-purple) and eosin (an acidic dye that stains basic/protein structures orange-red). This differential staining enables visualization of nuclear chromatin (blue-purple), cytoplasmic RNA (blue, as in polychromatophilic red cells), hemoglobin (pink-red), granules (neutrophilic --- lilac; eosinophilic --- orange-red; basophilic --- dark blue-black), and platelet cytoplasm (pale blue with purple granules). The thin area of the smear near the feathered edge, where red blood cells are evenly distributed without overlapping, is examined systematically. Modern laboratories may also use digital morphology systems (CellaVision, Medica) that scan and present cell images for review by a technologist or pathologist.

\section*{History}
The microscopic examination of blood dates to the 17th century when Antonie van Leeuwenhoek first observed red blood cells in 1674. However, the peripheral blood smear as a diagnostic tool emerged in the late 19th century. Paul Ehrlich, working in the 1870s--1880s, developed the first practical blood cell staining methods using aniline dyes, enabling differentiation of white blood cell types. He described neutrophils, eosinophils, basophils, and mast cells, and coined the term ``reticulocyte.'' Dmitri Romanowsky, a Russian physician, developed the first combined stain (methylene blue and eosin) in 1891, which uniquely demonstrated the purple color of nuclear chromatin and the red-pink color of hemoglobin --- the ``Romanowsky effect.'' James Homer Wright, an American pathologist, refined Romanowsky's technique in 1902 to produce the Wright stain, which remains the most widely used stain for peripheral blood smears in the United States. Gustav Giemsa modified the stain further in 1904. Throughout the 20th century, the peripheral blood smear was the central tool of hematology diagnosis. While automated analyzers have supplemented and largely replaced the manual differential for routine screening, the blood smear remains irreplaceable for morphological diagnosis and is still considered an essential skill in hematology training.

\section*{Why This Test Is Ordered}
\begin{itemize}
  \item Evaluation of abnormal automated CBC results flagged by the analyzer --- unexplained leukocytosis, leukopenia, thrombocytopenia, thrombocytosis, abnormal red cell indices, or abnormal cell populations
  \item Diagnosis and classification of anemia --- identifying specific morphologic abnormalities such as spherocytes (hereditary spherocytosis, autoimmune hemolytic anemia), sickle cells (sickle cell disease), target cells (thalassemia, liver disease, hemoglobinopathies), schistocytes (microangiopathic hemolytic anemia), acanthocytes (liver disease, abetalipoproteinemia), teardrop cells (myelofibrosis), and ovalocytes/elliptocytes (hereditary elliptocytosis)
  \item Diagnosis of hematologic malignancies --- identification of blast cells, Auer rods (acute myeloid leukemia), hairy cells (hairy cell leukemia), smudge cells (chronic lymphocytic leukemia), and abnormal lymphocytes
  \item Evaluation of unexplained jaundice, hemolysis, or splenomegaly
  \item Diagnosis and monitoring of infection --- toxic granulation, D\"{o}hle bodies, and vacuolization in neutrophils suggest bacterial infection; atypical lymphocytes (lymphocytosis) suggest viral infection (e.g., Epstein-Barr virus/infectious mononucleosis)
  \item Diagnosis of parasitic infections --- malaria (Plasmodium species), babesiosis, trypanosomiasis, leishmaniasis, and filariasis parasites
  \item Investigation of thrombocytopenia --- distinguishing between true thrombocytopenia and pseudothrombocytopenia (platelet clumping in EDTA), and examining platelet size (giant platelets in ITP, small platelets in Wiskott-Aldrich syndrome)
  \item Monitoring of chemotherapy or bone marrow recovery --- assessing for dysplastic features or regeneration
  \item Evaluation of unexplained constitutional symptoms --- fever, weight loss, night sweats, bone pain, or lymphadenopathy
\end{itemize}

\section*{Is This a Primary or Secondary Test}
The peripheral blood smear is a secondary (reflex) test, typically performed following abnormal automated CBC results or when specifically ordered based on clinical suspicion of a hematologic disorder.

\section*{How This Test Is Performed}
A venous blood sample is collected in a lavender-top (EDTA) tube from a peripheral vein, usually in the antecubital fossa. Alternatively, a fingerstick capillary blood sample can be used. A small drop of well-mixed blood is placed near the frosted end of a clean glass microscope slide. Using a second spreader slide held at a 30--45 degree angle, the drop of blood is spread in a single, smooth, rapid motion to create a thin film with a feathered edge. The slide is allowed to air-dry completely. It is then fixed in absolute methanol for 1--2 minutes and stained with Wright-Giemsa stain. Staining involves immersion in Wright stain for 2--3 minutes, addition of buffer, further staining for 5--10 minutes, rinsing with water, and air-drying. The stained slide is examined systematically under a microscope, first at low power (10$\times$) to assess overall quality, cellular distribution, and presence of platelet clumps or fibrin strands. Examination then proceeds to high dry (40$\times$) and oil immersion (100$\times$) objectives for detailed morphological assessment. The examiner performs a differential count (counting 100--200 white blood cells), assesses red blood cell morphology, and evaluates platelet number and morphology. The entire process, including slide preparation, staining, and expert review, typically takes 30--120 minutes depending on laboratory workflow and complexity.

\section*{What Preparation Is Needed}
No special preparation is required. Fasting is not necessary. The patient should inform the healthcare provider of any medications, particularly anticoagulants, antiplatelet agents, chemotherapy, immunosuppressants, or recent blood transfusion. Recent alcohol consumption should be noted, as it can produce vacuolization of myeloid precursors.

\section*{Contraindications and Precautions}
No absolute contraindications. Several pre-analytical factors critically affect smear quality and interpretation: the sample must be collected in EDTA (purple-top tube) and thoroughly mixed by gentle inversion; smears should be prepared within 2--4 hours of collection --- prolonged EDTA exposure causes degenerative changes including red cell sphering, nuclear pyknosis of white blood cells, and platelet swelling; clotted specimens are unacceptable; underfilling the EDTA tube (excess EDTA-to-blood ratio) causes red cell crenation and falsely elevated MCV; overfilling leads to inadequate anticoagulation and potential clotting. The smear must be of proper thickness and even distribution --- too thick prevents cell identification, too thin causes cell distortion. Cold agglutinins can cause red cell clumping and agglutination artifacts. Fingerstick samples must be free-flowing and not squeezed excessively, which can introduce tissue fluid and platelet clumps. Automated digital smear systems may be used as adjuncts but do not fully replace manual microscopic review for all abnormalities.

\section*{Risks}
Risks are those of routine venipuncture: mild pain, bruising, hematoma, lightheadedness, vasovagal syncope, and rarely, infection. Fingerstick carries minimal additional risk of local pain and infection. No additional risks are associated with the smear preparation and examination.

\section*{What Happens After the Test}
Smear review and reporting typically take 1--4 hours, depending on complexity and laboratory staffing. If significant abnormalities are identified (blast cells, Auer rods, severe plasmodial parasitemia, microangiopathic changes), the laboratory may notify the healthcare provider urgently (critical value alert). The pathologist's report describes the findings systematically: red blood cell morphology (anisocytosis, poikilocytosis, specific cell shapes, inclusions, polychromasia), white blood cell morphology (differential count, toxic changes, left shift, blasts, abnormal lymphocytes), and platelet assessment (number, size, morphology, clumping). Based on the smear findings, the pathologist or laboratory scientist provides an interpretive comment and recommendations for follow-up testing. The peripheral smear report is correlated with the automated CBC results, clinical history, and examination findings for final diagnostic interpretation.

\section*{Understanding Results}

\subsection*{Abnormal Red Blood Cell Morphology}
\begin{itemize}
  \item \textbf{Spherocytes:} Small, dense, spherical red cells lacking central pallor. Seen in hereditary spherocytosis (elevated MCHC, positive osmotic fragility test, negative direct Coombs) and autoimmune hemolytic anemia (positive direct Coombs test).
  \item \textbf{Schistocytes (fragmented cells):} Irregular, helmet-shaped, or triangular cell fragments indicating mechanical red cell fragmentation from microangiopathic hemolytic anemia. Associated with thrombotic thrombocytopenic purpura (TTP), hemolytic-uremic syndrome (HUS), disseminated intravascular coagulation (DIC), malignant hypertension, and prosthetic cardiac valves.
  \item \textbf{Sickle cells (drepanocytes):} Crescent-shaped, elongated cells with pointed ends. Pathognomonic for sickle cell disease (HbS). Confirmed by hemoglobin electrophoresis or solubility test.
  \item \textbf{Target cells (codocytes):} Cells with a central dark hemoglobin area surrounded by a clear ring and an outer dark rim, resembling a bull's-eye target. Associated with thalassemia, hemoglobin C disease, liver disease, post-splenectomy, and iron deficiency.
  \item \textbf{Teardrop cells (dacryocytes):} Pear-shaped red cells strongly suggestive of myelofibrosis with myeloid metaplasia, but also seen in thalassemia, megaloblastic anemia, and myelophthisic marrow infiltration.
  \item \textbf{Acanthocytes (spur cells):} Cells with irregular, asymmetrical, spiky surface projections. Seen in severe liver disease (spur cell anemia), abetalipoproteinemia, and post-splenectomy.
  \item \textbf{Echinocytes (burr cells):} Cells with regular, uniform, short spicules. Seen in uremia, pyruvate kinase deficiency, and as an artifact (EDTA storage, pH changes).
  \item \textbf{Elliptocytes and ovalocytes:} Elongated oval cells. Hereditary elliptocytosis (autosomal dominant) or megaloblastic anemia, iron deficiency, myelofibrosis.
  \item \textbf{Stomatocytes:} Cells with a slit-like central pallor (mouth-shaped). Hereditary stomatocytosis, liver disease, alcoholism.
  \item \textbf{Bite cells (degmacytes):} Cells with a semicircular ``bite'' removed. Characteristic of G6PD deficiency with oxidative hemolysis, associated with Heinz bodies (revealed by supravital stain).
  \item \textbf{Howell-Jolly bodies:} Small, round, dark purple nuclear remnants. Normally removed by the spleen; their presence indicates asplenia (surgical, functional, or congenital) or severe hemolysis with splenic overload.
  \item \textbf{Basophilic stippling:} Fine or coarse blue-purple granules representing aggregated ribosomal RNA. Seen in lead poisoning, thalassemia, sideroblastic anemia, and severe hemolysis.
  \item \textbf{Pappenheimer bodies:} Small, iron-containing basophilic granules. Confirmed by Prussian blue iron stain. Associated with sideroblastic anemia, thalassemia, and post-splenectomy.
  \item \textbf{Polychromasia:} Large, bluish-purple-staining red cells representing reticulocytes with residual RNA. Indicates active erythropoiesis (hemolysis, recovery from anemia).
  \item \textbf{Rouleaux formation:} Red cells stacked like coins. Associated with hypergammaglobulinemia (multiple myeloma, Waldenstr\"{o}m macroglobulinemia, chronic inflammatory conditions).
  \item \textbf{Autoagglutination:} Irregular clumping of red cells. Indicates cold agglutinin disease with IgM autoantibodies.
  \item \textbf{Intracellular organisms:} Malaria parasites (Plasmodium falciparum, vivax, ovale, malariae, knowlesi), Babesia species, or Bartonella bacilliformis within red blood cells.
\end{itemize}

\subsection*{Abnormal White Blood Cell Morphology}
\begin{itemize}
  \item \textbf{Toxic granulation:} Prominent, darkly staining primary granules in neutrophils. Indicates bacterial infection, inflammation, or toxic states (burns, chemotherapy).
  \item \textbf{D\"{o}hle bodies:} Pale blue-gray cytoplasmic inclusions representing aggregates of rough endoplasmic reticulum. Associated with bacterial infection, pregnancy, burns, and myelodysplasia.
  \item \textbf{Cytoplasmic vacuolization:} Vacuoles in neutrophils. Associated with severe bacterial infection/toxicity; also an artifact of prolonged EDTA storage.
  \item \textbf{Hypersegmented neutrophils:} Neutrophils with 6 or more nuclear lobes. Classical for megaloblastic anemia (B12 or folate deficiency).
  \item \textbf{Atypical (reactive) lymphocytes:} Large lymphocytes with abundant basophilic cytoplasm and irregular nuclei. Characteristic of infectious mononucleosis (Epstein-Barr virus), but also seen in cytomegalovirus, HIV, toxoplasmosis, and drug reactions.
  \item \textbf{Blast cells:} Large undifferentiated cells with fine chromatin, prominent nucleoli, and scant basophilic cytoplasm. Presence in peripheral blood indicates acute leukemia --- lymphoblasts (ALL) or myeloblasts (AML). Auer rods (linear cytoplasmic inclusions) are pathognomonic for AML (myeloid lineage).
  \item \textbf{Smudge (smear) cells:} Fragile, degenerated lymphocytes disrupted during smear preparation. Abundant smudge cells are characteristic of chronic lymphocytic leukemia (CLL).
  \item \textbf{Hairy cells:} Medium-sized lymphocytes with irregular cytoplasmic projections (``hairy'' borders). Diagnostic of hairy cell leukemia, confirmed by flow cytometry.
  \item \textbf{Pelger-Hu\"{e}t anomaly:} Neutrophils with bilobed (``spectacle'' or ``dumbbell'') nuclei without segmentation. Can be congenital (benign, autosomal dominant) or acquired (pseudo-Pelger-Hu\"{e}t, seen in myelodysplasia and acute myeloid leukemia).
  \item \textbf{Left shift (bands and metamyelocytes):} Increased immature myeloid precursors. Indicates acute bacterial infection, inflammation, or recovery from myelosuppression.
  \item \textbf{Leukoerythroblastic picture:} Presence of immature myeloid cells and nucleated red blood cells. Suggests bone marrow infiltration (myelofibrosis, metastatic carcinoma, lymphoma, granulomatous disease).
\end{itemize}

\subsection*{Abnormal Platelet Morphology}
\begin{itemize}
  \item \textbf{Giant platelets:} Platelets approaching or exceeding red blood cell size. Seen in immune thrombocytopenic purpura (ITP), Bernard-Soulier syndrome, and May-Hegglin anomaly.
  \item \textbf{Platelet clumping:} Usually a pre-analytical artifact from EDTA-induced platelet agglutination (pseudothrombocytopenia). Confirmed by re-collecting in citrate tube and observing corrected platelet count.
  \item \textbf{Platelet satellitism:} Platelets adhering to white blood cells (usually neutrophils), an EDTA artifact.
  \item \textbf{Hypogranular platelets:} Diminished platelet granules. Seen in myelodysplastic syndrome and acute myeloid leukemia.
\end{itemize}

\section*{Normal Results}
A normal peripheral blood smear demonstrates:
\begin{itemize}
  \item \textbf{Red blood cells:} Round, biconcave cells with a diameter of 6--8 $\mu$m and a central area of pallor occupying approximately 1/3 of the cell diameter (normocytic, normochromic). Minimal anisocytosis (size variation) and poikilocytosis (shape variation).
  \item \textbf{White blood cell differential:} Neutrophils 40--70\%, lymphocytes 20--40\%, monocytes 2--8\%, eosinophils 1--4\%, basophils 0--1\%. No immature myeloid cells (blasts, promyelocytes, myelocytes, metamyelocytes) and no nucleated red blood cells in the peripheral blood.
  \item \textbf{Platelets:} Small, anucleate, purple cytoplasmic fragments 1--4 $\mu$m in diameter, with 7--15 platelets per oil immersion field (corresponding to a platelet count of 150,000--400,000/$\mu$L). No clumping or giant forms.
  \item \textbf{No intracellular or extracellular organisms:} No malarial parasites, Babesia, trypanosomes, microfilariae, or bacteria.
\end{itemize}

\section*{Follow-up Tests}
\begin{itemize}
  \item \textbf{Complete blood count with automated differential:} Quantitative correlation with smear findings
  \item \textbf{Reticulocyte count:} To quantify erythropoietic response when polychromasia is present on smear
  \item \textbf{Iron studies (serum iron, ferritin, TIBC):} When microcytic hypochromic cells are noted
  \item \textbf{Vitamin B12 and serum folate levels:} When hypersegmented neutrophils and macro-ovalocytes suggest megaloblastic anemia
  \item \textbf{Direct antiglobulin test (Coombs test):} When spherocytes and polychromasia suggest autoimmune hemolytic anemia
  \item \textbf{Hemoglobin electrophoresis:} When sickle cells, target cells, or other abnormal erythrocyte morphology suggests hemoglobinopathy
  \item \textbf{Flow cytometry/immunophenotyping:} When blast cells or abnormal lymphocytes are identified, to determine lineage and clonality
  \item \textbf{Bone marrow aspiration and biopsy:} When blast cells, leukoerythroblastic picture, or unexplained cytopenias are present
  \item \textbf{Cytogenetic analysis (karyotype, FISH):} When acute leukemia or myelodysplasia is suspected
  \item \textbf{Malaria rapid diagnostic test (RDT) and thick blood smear:} When malaria parasites are suspected or seen on thin smear (species identification and parasitemia quantification)
  \item \textbf{G6PD enzyme assay:} When bite cells, Heinz bodies, or acute hemolysis are present
  \item \textbf{Osmotic fragility test:} When spherocytes are present with a negative direct Coombs test
  \item \textbf{ADAMTS13 activity assay:} When schistocytes and thrombocytopenia suggest thrombotic thrombocytopenic purpura
  \item \textbf{Coagulation panel (PT, PTT, fibrinogen, D-dimer):} When schistocytes and thrombocytopenia suggest DIC
\end{itemize}

\section*{References}
\begin{enumerate}
  \item Bain BJ. Blood Cells: A Practical Guide. 6th ed. Wiley Blackwell; 2022.
  \item Bain BJ, Bates I, Laffan MA. Dacie and Lewis Practical Haematology. 12th ed. Elsevier; 2017.
  \item Kaushansky K, Lichtman MA, Prchal JT, et al. Williams Hematology. 10th ed. McGraw-Hill; 2021.
  \item Tierney LM, Saint S, Drazen JM. CURRENT Medical Diagnosis \& Treatment. McGraw-Hill; 2023.
  \item Wallach J. Interpretation of Diagnostic Tests. 11th ed. Wolters Kluwer; 2022.
  \item Fischbach FT, Dunning MB. A Manual of Laboratory and Diagnostic Tests. 11th ed. Wolters Kluwer; 2023.
  \item Palmer L, Briggs C, McFadden S, et al. ICSH recommendations for the standardization of nomenclature and grading of peripheral blood cell morphological features. Int J Lab Hematol. 2015;37(3):287--303.
  \item Rodak BF, Carr JH. Clinical Hematology Atlas. 5th ed. Elsevier; 2017.
\end{enumerate}

\clearpage
